\section{Introduction}

Mangrove canopy height describes an important dimension of forest structure that cannot be recovered from an extent map alone. Differences in height can inform interpretation of stand development and structural heterogeneity, and height is relevant to biomass estimation when combined with appropriate ecological information and allometric relationships. Global mangrove research has linked canopy-height variation to environmental gradients and examined its relationship with carbon stocks \citep{simard2019}. For a regional mapping study, the immediate objective is therefore to estimate forest structure consistently across space; translating those estimates into biomass or carbon requires additional validation.

Earth observation offers an opportunity to complement local structural measurements with spatially extensive imagery. GEDI provides observations of vertical forest structure from spaceborne waveform lidar \citep{dubayah2020}, but its sampling pattern leaves gaps between observations. Sentinel imagery supplies complementary spatial coverage. Learning a mapping from these images to sparse height references can produce predictions at locations without lidar measurements. Existing work already demonstrates the feasibility of learning dense canopy-height maps from sparse GEDI supervision and Sentinel-2 imagery \citep{lang2023}. The present study extends this established formulation to a controlled regional comparison of optical--SAR fusion strategies.

The unresolved question addressed here concerns the choice of optical--SAR fusion strategy in a regional mangrove setting. Sentinel-2 spectral bands and vegetation indices provide optical information, while Sentinel-1 contributes radar backscatter measurements. Combining the modalities may improve prediction, but improvement cannot be assumed simply because more channels are available. Input-level concatenation requires a single encoder to learn from both modalities. Separate encoders allow modality-specific representations before feature fusion, but also alter model capacity. A controlled comparison is needed to distinguish the empirical value of additional SAR inputs from the value of a different architecture.

Cà Mau provides the regional setting for this investigation. The study uses the mangrove extent represented by the supplied mask and Sentinel composites for 2020--2021. Coastal landscape boundaries make the distinction between reference support and map resolution especially relevant: a GEDI footprint is wider than a 10 m prediction cell, and a footprint near a forest boundary can represent mixed conditions. Specific local descriptions of forest composition, management, tidal conditions, and restoration history require regional evidence and are left for verification \TODO{SOURCE NEEDED}.

Evaluation is equally central to the research question. Nearby image patches and lidar observations can share spatially structured information. Consequently, randomly separating observations does not directly test prediction in a geographically withheld region. We assign three contiguous regions along easting before selecting patches and exclude buffered boundary areas. This design measures transfer within the study domain across a fixed spatial partition. It does not establish transfer to other provinces or independence at all spatial scales.

We make three application-focused contributions. First, we implement a regional sparse-to-dense regression workflow in which GEDI RH98 remains sparse and supervises predictions only at valid reference cells; deployment uses Sentinel inputs alone. Second, we compare optical-only, SAR-only, early-fusion, and dual-encoder models within a shared experimental protocol to quantify the observed contribution of modality and fusion design. Third, we report geographically separated evaluation together with the actual reference counts and limitations, distinguishing prediction-grid spacing from lidar support and sampled accuracy from map-wide accuracy.
